%%%%%%%%%%%%%%%%%%%%%%%%%%%%%%%%%%%%%%%%%%%%%%%%%%%%%%%%%%%%%%%%%%%%%%%%%%%%%%%
% ARTICLE 16: Competition Teams
%%%%%%%%%%%%%%%%%%%%%%%%%%%%%%%%%%%%%%%%%%%%%%%%%%%%%%%%%%%%%%%%%%%%%%%%%%%%%%%

\Article{Competition Teams}

\begin{enumerate}
	\item A Competition Team is a group of RIT students organized to prepare for and compete 
            in the specific external cybersecurity competition for which the team is formed, 
            representing RIT as the university’s official team for that competition.
	\item E-Board reserves the right to amend 16 RITSEC without express club approval and/or 
            vote; however, the Coaches and Cybersecurity Department must be consulted before any 
            changes are finalized, as they serve as the official university representatives to the 
            competition organizers.
      \item Captains of the Competition Teams that are sponsored by the RIT cybersecurity department 
            will be given A-Board positions.
      \item The E-Board may, by majority vote, recommend to the Coaches and Department that a 
            Competition Team be disbanded, a captain be removed, or other intervention occur if a 
            violation of the RITSEC Code of Conduct (defined in Article 9, Section A) is believed 
            to have taken place. Final authority for such actions rests solely with the Coach(es) 
            as the official university representatives.
      \item The club may, by majority vote, provide additional input or recommendations to the Coaches 
            regarding any action proposed under Section IV.
      \item Competition Team Captains, after discussion with the Coach(es), may relay to RITSEC any 
            sponsorship opportunities and major purchase plans for the team.
      \item All Competition Team members are considered members of RITSEC and are therefore subject 
            to Article 9.

\end{enumerate}

% SECTION A: Funding
\Section{Funding}

\begin{enumerate}
	\item Competition Teams will be able to propose a budget to RITSEC, contingent upon the availability of funds within the RITSEC budget. These funds may support competition fees, training materials, competition travel expenses, and necessary licensing for technologies used in competition. Additionally, Teams may seek funding to cover certification costs for technologies relevant to competition scenarios. Any budget received from RITSEC will be supplemental to existing funding from the department.
\end{enumerate}


% SECTION D: CCDC
\Section{Collegiate Cyber Defense Competition Team}

% SUBSECTION 1: NAME
\Subsection{Name}

\begin{enumerate}
	\item Collegiate Cyber Defense Competition (CCDC) Team
\end{enumerate}

% SUBSECTION 2: PURPOSE
\Subsection{Purpose}

\begin{enumerate}
	\item To prepare for and compete in the National Collegiate Cyber Defense Competition.
	\item To compete in other cyber defense competitions as a means of refining skills and enhancing readiness for CCDC.
      \item To provide educational opportunities through information sessions and a tryout process for students 
            interested in learning more about cyber defense, incident response, and blue teaming, as well as 
            those considering joining the team.
\end{enumerate}

% SUBSECTION 3: TARGET AUDIENCE
\Subsection{Target Audience}

\begin{enumerate}
	\item Official members of the RIT CCDC team, as selected by the coaches via the official tryout process.
      \item Students interested in trying out for the team who both are passionate about the subject area 
            of the competition and have specialized experience which fits the needs of the team.
\end{enumerate}

% SUBSECTION 4: REQUIREMENTS
\Subsection{Requirements}

\begin{enumerate}
	\item The Team Captain is responsible for proposing a successor to the Coaches upon graduation, 
            departure from the team, or resignation from the Captain role; the Coaches shall make the 
            final appointment as required by CCDC rules.
	\item The team shall adhere to all official national rules for the Collegiate Cyber Defense 
            Competition, as well as any local variations specific to the Northeast region.
	\item The Coach(es), as designated by the Department, shall serve as the university’s official 
            representatives to competition organizers and will coordinate travel itineraries, 
            accommodations, and other competition logistics. The Coach(es) will accompany the team 
            when required by university or competition policy or when they determine it is necessary. 
            The Coach(es) shall schedule and facilitate a mandatory pre-travel meeting for participating 
            team members, while any additional Coach attendance at team meetings will be determined by the 
            Coach(es) based on team needs, applicable rules, and availability.
      \item Membership is contingent upon being a student in good standing as defined by RIT’s Undergraduate 
            Policy or Graduate Policy (RIT D05.1 § II.A-B). Additional factors include, but are not limited to, 
            participation in tryouts and evaluations conducted by current team members. Final decisions regarding 
            membership will be made by the Coach(es), taking into account recommendations from the team.
      \item The team will be able to utilize RITSEC resources as necessary for team activities, including practices 
            and the CCDC “Mini Competition” as part of tryouts.
\end{enumerate}
% SECTION D: CCDC
\Section{Collegiate Penetration Testing Competition Team}

% SUBSECTION 1: NAME
\Subsection{Name}

\begin{enumerate}
	\item Collegiate Penetration Testing Competition (CPTC) Team
\end{enumerate}

% SUBSECTION 2: PURPOSE
\Subsection{Purpose}

\begin{enumerate}
	\item To train and prepare to qualify for and compete in any CPTC Region Competition.
      \item To provide educational opportunities through information sessions and tryout sessions 
            for students interested in learning more about cyber defense, incident response, and blue 
            teaming, as well as those considering joining the team.
\end{enumerate}

% SUBSECTION 3: TARGET AUDIENCE
\Subsection{Target Audience}

\begin{enumerate}
	\item Official members of the RIT CPTC team.
      \item Students interested in trying out for the team who both are passionate about the subject area 
            of the competition and have specialized experience which fits the needs of the team.
\end{enumerate}

% SUBSECTION 4: REQUIREMENTS
\Subsection{Requirements}

\begin{enumerate}
\item To foster interest in the competition and improve development of skills in the competition, the team 
      will host three practices open to all RITSEC members during the off-season. Membership is at the complete 
      discretion of the captain and the team’s faculty advisor.
\item The team shall adhere to all official national rules for the Collegiate Penetration Testing Competition.
\item The coach, as designated by the department/team, will be leading the organization of the travel itinerary, 
      accommodation, and other logistics related to competitions, as well as accompany the team on travel. The 
      coach is expected to attend at least three team meetings throughout the season, including an introductory 
      meeting, a trip planning session, and a post-competition debrief.
\item Membership is contingent upon being a student in good standing as defined by RIT’s Academic Policy. Additional 
      factors include, but are not limited to, participation in tryouts and evaluations conducted by current team 
      members. Final decisions regarding membership will be made by the coach, taking into account recommendations 
      from the team.

\end{enumerate}